\documentclass[tikz,border=5pt]{standalone}
\usepackage{esint}
\tikzset{
    pi/.pic={\node[rotate=#1]{$\pi$};},
    pipic/.pic={
        \pic[orange] at (0,-.1) {pi=0};
        \pic[red] at (0,.1) {pi=180};
        \pic[magenta] at (-.16,0) {pi=90};
        \pic[teal] at (.16,0) {pi=270};
    }
}
\newcommand*\ljhknot[1]{%
    \tikz{\foreach \x in {0,...,#1}{
        \foreach \y in {0,...,#1}{
            \pic at ({(\x-\y)*.25},{(\x+\y)*.27}) {pipic};
        }
        \node[red,yshift={-.4cm*(\x+1)}] {$\iint$};
    }}
}
\begin{document}

% \begin{tikzpicture}[
%     pi/.pic={\node[rotate=#1]{$\pi$};},
%     pipic/.pic={
%         \pic[orange] at (0,-.1) {pi=0};
%         \pic[red] at (0,.1) {pi=180};
%         \pic[magenta] at (-.16,0) {pi=90};
%         \pic[teal] at (.16,0) {pi=270};
%     }
% ]
% \pic {pipic};
% \pic at (.25,.27) {pipic};
% \def\NN{2}
% \foreach \x in {0,...,\NN}{
%     \foreach \y in {0,...,\NN}{
%         \pic at ({(\x-\y)*.25},{(\x+\y)*.27}) {pipic};
%     }
%     \node[red,yshift={-.4cm*(\x+1)}] {$\iint$};
% }
% \end{tikzpicture}

\ljhknot{2}

\ljhknot{8}

\end{document}